\poczatekwykladu{7}
\section{Dystrybuanty miar}
\label{sec:dystrybuanty-miar}

Konstrukcja Carathéodory'ego dała nam przestrzeń z miarą
$(\R,\mathcal{L},\lambda)$, przy czym
$\mathcal{B}(\R)\subseteq\mathcal{L}$. Obcięcie $\lambda$ do
$\mathcal{B}(\R)$ również jest miarą. Rozważymy teraz inne miary
na prostej oraz ich opis za pomocą funkcji rzeczywistych.

\begin{definicja}[Dystrybuanta miary]
\label{def:dystrybuanta-miary}
Niech $\mu$ będzie miarą na $\mathcal{B}(\R)$, skończoną na każdym
ograniczonym przedziale. Jej dystrybuantą nazywamy funkcję
$F_\mu\colon\R\to\R$ określoną wzorem
\begin{equation}
  F_\mu(t)=
  \begin{cases}
    \mu([0,t]), & t\geq0,\\
    -\mu((t,0)), & t<0.
  \end{cases}
  \label{eq:dystrybuanta-miary}
\end{equation}
\end{definicja}

\begin{fakt}[Własności dystrybuanty]
\label{fakt:dystrybuanta-wlasnosci}
Dystrybuanta $F_\mu$ ma następujące własności:
\begin{enumerate}
  \item jest niemalejąca;
  \item jest prawostronnie ciągła, czyli dla każdego $t_0\in\R$
    \[
      \lim_{t\to t_0^+}F_\mu(t)=F_\mu(t_0);
    \]
  \item $\displaystyle\lim_{t\to0^-}F_\mu(t)=0$.
\end{enumerate}
Ponadto dla $a<b$ zachodzi
\begin{equation}
  \mu((a,b])=F_\mu(b)-F_\mu(a).
  \label{eq:miara-przez-dystrybuante}
\end{equation}
\end{fakt}

\begin{proof}
Równość~\eqref{eq:miara-przez-dystrybuante} wynika z addytywności:
dla $0\leq a<b$ odejmujemy miary $[0,b]$ i $[0,a]$, dla $a<b<0$
rozkładamy $(a,0)=(a,b]\cup(b,0)$, a dla $a<0\leq b$ używamy
$(a,b]=(a,0)\cup[0,b]$. Nieujemność miary daje monotoniczność $F_\mu$.
Jeżeli $t_n\downarrow t_0$ i $t_n>t_0$, to
$(t_0,t_n]\downarrow\varnothing$, zatem
$F_\mu(t_n)-F_\mu(t_0)=\mu((t_0,t_n])\to0$.
Stosujemy tu ciągłość z góry do miary obciętej do ograniczonego
przedziału. Podobnie dla $t_n\uparrow0$, $t_n<0$, mamy
$(t_n,0)\downarrow\varnothing$ i $F_\mu(t_n)\to0$.
\end{proof}

\begin{twierdzenie}[Odpowiedniość miar i dystrybuant]
\label{tw:miary-dystrybuanty}
Niech $\mathcal{D}$ oznacza rodzinę wszystkich funkcji
$F\colon\R\to\R$, które są niemalejące, prawostronnie ciągłe
i spełniają $\lim_{t\to0^-}F(t)=0$.
Odwzorowanie $\mu\mapsto F_\mu$ jest bijekcją między rodziną miar
na $\mathcal{B}(\R)$, skończonych na przedziałach ograniczonych,
a rodziną $\mathcal{D}$. Miara odpowiadająca funkcji $F$ spełnia
\[
  \mu((a,b])=F(b)-F(a)\qquad(a<b).
\]
\end{twierdzenie}

\begin{proof}[Idea dowodu]
Dla danej funkcji $F\in\mathcal{D}$ przypisujemy przedziałom wartości
$m((a,b])=F(b)-F(a)$. Następnie definiujemy miarę zewnętrzną przez
pokrycia takimi przedziałami i stosujemy konstrukcję Carathéodory'ego.
Otrzymana miara borelowska ma wskazane wartości na przedziałach.
Jest przez nie wyznaczona jednoznacznie. Normalizacja
$\lim_{t\to0^-}F(t)=0$ ustala stałą addytywną w dystrybuancie.
\end{proof}

\begin{samepage}
\begin{przyklad}[Miara Lebesgue'a]
\label{prz:dystrybuanta-lebesgue}
Dla miary Lebesgue'a obciętej do $\mathcal{B}(\R)$ mamy
\[
  F_\lambda(t)=
  \begin{cases}
    t, & t\geq0,\\
    -(-t)=t, & t<0.
  \end{cases}
\]
Zatem $F_\lambda(t)=t$ na całej prostej.
\end{przyklad}

\begin{przyklad}[Miara Diraca w zerze]
\label{prz:dystrybuanta-diraca}
Dla miary Diraca $\delta_0$ dystrybuanta ma postać
\[
  F_{\delta_0}(t)=
  \begin{cases}
    0, & t<0,\\
    1, & t\geq0.
  \end{cases}
\]
W zerze występuje skok, a wartość w tym punkcie jest równa $1$,
zgodnie z prawostronną ciągłością. Wykresy obu dystrybuant
przedstawia rysunek~\ref{fig:dystrybuanty-lebesgue-dirac}.
\end{przyklad}
\end{samepage}

\begin{figure}[htbp]
  \centering
  \includegraphics[width=0.95\textwidth]{dystrybuanty-lebesgue-dirac.pdf}
  \caption{Dystrybuanty miary Lebesgue'a i miary Diraca w zerze.}
  \label{fig:dystrybuanty-lebesgue-dirac}
\end{figure}

\newpage
\begin{przyklad}[Miara skupiona w trzech punktach]
\label{prz:dystrybuanta-trzy-punkty}
Rozważmy miarę z tablicy
\[
  \mu=\frac14\delta_{-1}+\frac12\delta_0+\frac14\delta_1.
\]
Dla $t<0$ przedział $(t,0)$ nie zawiera punktu $0$.
Zawiera punkt $-1$ dokładnie wtedy, gdy $t<-1$.
Dla $t\geq0$ przedział $[0,t]$ zawiera punkt $0$, a punkt $1$
dokładnie wtedy, gdy $t\geq1$. Stąd
\[
  F_\mu(t)=
  \begin{cases}
    -\frac14, & t<-1,\\
    0, & -1\leq t<0,\\
    \frac12, & 0\leq t<1,\\
    \frac34, & t\geq1.
  \end{cases}
\]
Wykres przedstawia rysunek~\ref{fig:dystrybuanty-atomowa-szkic} po lewej.
Normalizacja z definicji~\ref{def:dystrybuanta-miary} daje tu
ujemne wartości dla $t<-1$.
\end{przyklad}

\begin{przyklad}[Dystrybuanta miary na zbiorze Cantora]
\label{prz:dystrybuanta-cantora}
Trzeci wykres z tablicy przedstawia dystrybuantę miary skupionej
na zbiorze Cantora w $[0,1]$.
Na rysunku~\ref{fig:dystrybuanty-atomowa-szkic} po prawej odtworzono
jej schematyczny wykres. Widoczne są poziome fragmenty,
w szczególności między $1/3$ a $2/3$, na wysokości $1/2$.
\end{przyklad}

\begin{figure}[htbp]
  \centering
  \includegraphics[width=0.98\textwidth]{dystrybuanty-atomowa-szkic.pdf}
  \caption{Dystrybuanta miary z trzema atomami oraz dystrybuanta miary na zbiorze Cantora (schemat).}
  \label{fig:dystrybuanty-atomowa-szkic}
\end{figure}
